% =====================================================================
%  Chapter 21 review -- the 14-station revision contract.
%  Revision-only: tables, one-liners, pointers back via \secref/\thmref/
%  \propref/\defref/\coderef/\figref/\probref.  No re-derivations.
%  Summit: CF 840D Destiny (2500) -- the minimal value that occurs in a
%  STATIC range [l,r] STRICTLY MORE THAN (r-l+1)/k times (2<=k<=5), else -1.
%  n,q <= 3e5, 1 <= a_i <= n.  Two engines from this chapter combine:
%    * the merge-sort tree of 21.1 counts occurrences of a candidate value
%      in [l,r] in O(log^2 n) by bisection, so a random sample of O(1)
%      positions of [l,r] finds every ">len/k" value with high probability
%      (at most k-1 such values exist) -- the canonical Destiny editorial;
%    * the chairman/prefix-root engine of 21.3 gives the same frequency
%      deterministically (count of value v in [l,r] = two-root difference).
%  Both official samples hand-replayed:
%    4 2 / 1 1 2 2 / (1 3 2)->1 / (1 4 2)->-1
%    5 3 / 1 2 1 3 2 / (2 5 3)->2 / (1 2 3)->1 / (5 5 2)->2
%  Exports prob:c19-summit, lst:c19-kit, lst:c19-summit.
%  Mixed set is EXACTLY 10, all live-verified, un-owned, sorted ASCENDING by
%  effective difficulty; each hint names its reduction and cites a c19 listing.
%  CF 459D / 522D / 1000F were drafted here but are OWNED by Ch.18 (c14) --
%  they re-enter ONLY as prose re-treats in Related Topics, never re-linked.
% =====================================================================

\chaptersection{Chapter 21 Summary}

One sentence: \emph{when the data never changes but the questions multiply -- \emph{how
many values \(\le x\) in this range}, \emph{the \(k\)-th smallest here}, \emph{what was
stored at that past version}, \emph{the minimum of that block} -- one mutable aggregate
is the wrong object, and the fix is to store more per node (a sorted list), to store more
nodes (path-copied versions), or to store more lengths (a sparse table).} Chapter~19 could
fold one aggregate per interval and edit a point; Chapter~20 could rewrite a whole interval
with a lazy tag. This chapter is for the regimes those two cannot serve: a \emph{value} axis
sitting inside every \emph{index} node (the merge-sort tree), history you refuse to destroy
(persistence by path copying), prefixes of time as subtractable roots (chairman trees), the
same roots planted on a tree (four-root paths), and pure static \(\min/\max\) read in
constant time (the sparse table). Each row is a pointer, not a proof.

\begin{center}
\renewcommand{\arraystretch}{1.25}
\begin{tabular}{@{}p{0.24\linewidth} p{0.48\linewidth} p{0.20\linewidth}@{}}
\toprule
\textbf{Result} & \textbf{Statement} & \textbf{Where} \\
\midrule
Merge-sort tree & Each node stores its interval's \emph{sorted list}; count \(\le x\) in
\([l,r]\) is one bisection per canonical node; \(k\)-th is that count binary-searched on
value. & \defref{c19-mst}, \thmref{c19-mstcount} \\
Path copying & A new version shares every untouched subtree and copies only the
\(\bigO{\log n}\) root-to-leaf path; a node pool sizes the whole structure. & \defref{c19-pathcopy}, \thmref{c19-memory} \\
Persistent array / tree & Version roots index one shared structure; any past root is read
in \(\bigO{\log n}\) with no rebuild. & \coderef{c19-parray}, \coderef{c19-pseg} \\
Chairman prefixes & \(\mathit{root}_i\) = value-axis frequency tree after \(A[1..i]\);
count \(\le x\) in \([l,r]\) is \(\mathit{root}_r\!-\!\mathit{root}_{l-1}\), \(k\)-th is a
joint two-root descent. & \defref{c19-roots}, \thmref{c19-kth} \\
Roots on a tree & One root per DFS node holds the root-to-\(v\) path multiset; a path
\(u\text{--}v\) is a \emph{four}-root combination. & \defref{c19-rootpath}, \propref{c19-fourroots} \\
Sparse table & \(\mathit{st}[j][i]\) merges the \(2^j\) block from \(i\); \(\bigO{n\log n}\)
build, \(\bigO{1}\) query by two \emph{overlapping} blocks -- legal only for idempotent
merges. & \defref{c19-spacetable}, \thmref{c19-oone} \\
\bottomrule
\end{tabular}
\end{center}

\medskip
{\small\textbf{Algorithms from this chapter.}}
\begin{center}
\renewcommand{\arraystretch}{1.2}
\begin{tabular}{@{}p{0.34\linewidth} l l >{\raggedright\arraybackslash}p{0.22\linewidth}@{}}
\toprule
\textbf{Algorithm} & \textbf{Prep} & \textbf{Per op} & \textbf{Code} \\
\midrule
Merge-sort tree build (merge sorted lists) & \(\bigO{n\log n}\) & -- & \coderef{c19-mst-build} \\
Range count \(\le x\) (bisection per node) & \(\bigO{n\log n}\) & \(\bigO{\log^2 n}\) & \coderef{c19-mst-count} \\
SPOJ KQUERY (count \(>k\) in a range) & \(\bigO{n\log n}\) & \(\bigO{\log^2 n}\) & \coderef{c19-kquery-soln} \\
Persistent array (path-copied roots) & \(\bigO{n}\) & \(\bigO{\log n}\) & \coderef{c19-parray} \\
Persistent point-update / range-sum tree & pool & \(\bigO{\log n}\) & \coderef{c19-pseg} \\
CF 707D Persistent Bookcase (version DFS) & \(\bigO{qm}\) & \(\bigO{1}\) & \coderef{c19-707-soln} \\
Chairman build (one insert per element) & \(\bigO{n\log V}\) & \(\bigO{\log V}\) & \coderef{c19-chairman} \\
\(k\)-th by joint two-root descent & \(\bigO{n\log V}\) & \(\bigO{\log V}\) & \coderef{c19-kth} \\
SPOJ MKTHNUM (K-th Number) & \(\bigO{n\log n}\) & \(\bigO{\log n}\) & \coderef{c19-mkthnum-soln} \\
Persistent root per DFS node & \(\bigO{N\log V}\) & \(\bigO{\log V}\) & \coderef{c19-persist-tree} \\
Four-root \(k\)-th on a path & \(\bigO{N\log V}\) & \(\bigO{\log V}\) & \coderef{c19-pathkth} \\
SPOJ COT (Count on a tree) & \(\bigO{N\log N}\) & \(\bigO{\log N}\) & \coderef{c19-cot-soln} \\
Sparse table (idempotent \(O(1)\) RMQ) & \(\bigO{n\log n}\) & \(\bigO{1}\) & \coderef{c19-sparse} \\
Production persistent kit & pool & \(\bigO{\log n}\) & \coderef{c19-kit} \\
CF 840D Destiny (Summit) & \(\bigO{n\log n}\) & \(\bigO{\log^2 n}\) & \coderef{c19-summit} \\
\bottomrule
\end{tabular}
\end{center}

\chaptersection{Master Decision Flowchart}

Post-study ladder: walk top to bottom, stop at the first branch that matches. It decides
\emph{which static / persistent engine the query is begging for} -- then hands off what this
chapter does \emph{not} own: live range \emph{updates} (lazy, Ch.~20 \(\chref{c15}\)),
path updates on a changing tree (HLD, Ch.~22 \(\chref{c16}\)), non-idempotent range
combines that also mutate (beats / Li Chao, Ch.~23 \(\chref{c29}\)), and everything offline
(Mo's algorithm, Ch.~27 \(\chref{c21}\)).

\begin{lstlisting}[language={}, basicstyle=\ttfamily\small, numbers=none, frame=single]
  The array/tree is FIXED after build; you only READ it, many times.  Which engine?
  |
  |-- query is about VALUES in a range: "how many of a[l..r] are <= x"?
  |     +-- MERGE-SORT TREE (21.1, def:c19-mst).  Each node holds its interval's
  |        sorted list; count with a bisection per canonical node -> O(log^2 n).
  |        lst:c19-mst-build / lst:c19-mst-count / lst:c19-kquery-soln.
  |
  |-- query is the k-TH SMALLEST of a[l..r] (one fixed array, many queries)?
  |     +-- CHAIRMAN / PREFIX-VERSION tree (21.3, def:c19-roots).  root_i is the
  |        value histogram of a[1..i]; k-th is a joint descent of root_(l-1) and
  |        root_r in O(log V).  O(log^3 n) merge-sort + value-bisect is the
  |        fallback when you already have 21.1 built.  lst:c19-kth.
  |
  |-- values live on a TREE and the query is "k-th / count on path u--v"?
  |     +-- PERSISTENT ROOTS ON A TREE (21.4, prop:c19-fourroots).  One frequency
  |        root per DFS node; combine FOUR roots (u, v, lca, parent(lca)).  lca is
  |        Ch.11's (c11).  lst:c19-persist-tree / lst:c19-pathkth.
  |
  |-- pure static MIN / MAX / GCD of a[l..r], ZERO updates, read >> n times?
  |     +-- SPARSE TABLE (21.5, def:c19-spacetable).  O(1) by two OVERLAPPING
  |        blocks -- idempotent merges only.  NOT sum (prefix sums, Ch.18 c14).
  |        lst:c19-sparse.
  |
  |-- the data DOES change, but queries ask about a PAST version?
  |     +-- PERSISTENCE BY PATH COPYING (21.2, def:c19-pathcopy).  Copy the
  |        root-to-leaf path, share the rest; keep one root per version.  The
  |        node pool, not 4n, sizes it (thm:c19-memory).  lst:c19-pseg.
  |
  |-- updates are on the CURRENT state and need a whole-range rewrite?
  |     +-- NOT this chapter: lazy tags (Ch.20, c15).  A tag lives in no single
  |        version, so persistence and lazy refuse to share -- secref:c15-strains.
  |
  |-- need partial persistence by UNDO (rollback) instead of full roots?
  |     +-- forward-cite Ch.31 (c17) -- rollback DSU, not path copying.
\end{lstlisting}

\chaptersection{Common Mistakes Checklist}

\begin{itemize}[nosep, leftmargin=1.4em]
  \item \textbf{Path-copying only part of the path, or forgetting the leaf.} A persistent
        update must clone EVERY node on the root-to-leaf chain \emph{including} the leaf it
        rewrites (\(\defref{c19-pathcopy}\)); sharing the leaf you mutate corrupts the old
        version -- the persistence sin. \(\coderef{c19-pseg}\).
  \item \textbf{Mutating a shared child pointer.} In a persistent tree a child node may be
        reachable from several roots; writing through it retroactively edits the past. Copy
        first, then write. \(\secref{c19-copy}\).
  \item \textbf{Node-pool budget miscalculation \(\to\) MLE.} The cost is
        \(n + q\log n\) (or \(N\log V\) on a tree), not \(4n\); size the pool for the number
        of \emph{copies} or you run out (\(\thmref{c19-memory}\), \(\coderef{c19-pseg}\)).
  \item \textbf{Two-root subtraction off by one: use \(\mathit{root}_{l-1}\).} The range
        \([l,r]\) is \(\mathit{root}_r - \mathit{root}_{l-1}\), never \(\mathit{root}_l\) --
        index \(l\) belongs to the range (\(\thmref{c19-kth}\), \(\coderef{c19-kth}\)).
  \item \textbf{Four-root path double-counts the LCA.} The extra root is
        \(\mathit{root}_{\mathrm{parent}(\ell)}\), not \(\mathit{root}_\ell\); subtracting
        the LCA itself deletes it from the path (\(\propref{c19-fourroots}\),
        \(\coderef{c19-pathkth}\)).
  \item \textbf{Sparse table overlap is only valid for idempotent merges.} Sum, product, or
        popcount double-count the intersection; the \(O(1)\) two-block trick is licensed by
        \(\mathcal{M}(x,x)=x\) alone (\(\thmref{c19-oone}\), \(\coderef{c19-sparse}\)).
  \item \textbf{Log-of-one range edge.} A length-\(1\) query has \(j=\lfloor\log_2 1\rfloor=0\);
        guard \(\mathit{lg}[1]=0\) or the shift underflows (\(\coderef{c19-sparse}\)).
  \item \textbf{Bisection bounds when turning count into \(k\)-th.} The merge-sort-tree
        \(k\)-th binary-searches the \emph{value} axis; a loose upper bound or an
        off-by-one in \(\le\) vs \(<\) misses the smallest qualifying value
        (\(\secref{c19-mst}\)).
  \item \textbf{Coordinate compression forgotten.} Chairman trees index the \emph{value}
        axis; \(|a_i|\) up to \(10^9\) needs ranks \([1..\text{distinct}]\) before any insert
        (\(\defref{c19-roots}\), \(\coderef{c19-chairman}\)).
\end{itemize}

\chaptersection{WA Clinic}

Verdicts are symptoms; the causes are the short list this chapter drilled. Diagnose
top-down -- every leaf names a listing.

\begin{figure}[H]
\centering
\begin{tikzpicture}[diagram,
    v/.style={boxnode, fill=primaryblue!12, draw=primaryblue, text width=22mm, align=center},
    b/.style={boxnode, fill=difforange!12, draw=difforange, text width=12mm},
    l/.style={boxnode, fill=successgreen!10, draw=successgreen, text width=48mm,
              font=\scriptsize, align=left}]
  \node[v] (verdict) {\textbf{Verdict on the judge}};
  \node[b, below left=8mm and 5mm of verdict] (wa) {\textbf{WA}};
  \node[b, below right=8mm and 5mm of verdict] (mle) {\textbf{TLE / MLE}};
  \node[l, below right=1mm and 2mm of wa.south] (wa1)
    {a two-root range read used \(\mathit{root}_l\) instead of \(\mathit{root}_{l-1}\), or a
     path combined \(\mathit{root}_\ell\) rather than \(\mathit{root}_{\mathrm{parent}(\ell)}\)
     (\(\coderef{c19-kth}\), \(\coderef{c19-pathkth}\))};
  \node[l, below left=1mm and 2mm of wa.south] (wa2)
    {a sparse-table query on a NON-idempotent merge (sum) double-counted the overlap, or the
     two blocks failed to cover \([l,r]\) (\(\thmref{c19-oone}\), \(\coderef{c19-sparse}\))};
  \node[l, below right=1mm and 2mm of mle.south] (mle1)
    {node pool sized \(4n\) not \(n+q\log n\); every \code{newnode} leaked a shared child
     instead of reusing it (\(\thmref{c19-memory}\), \(\coderef{c19-pseg}\))};
  \node[l, below left=1mm and 2mm of mle.south] (mle2)
    {merge-sort tree queried by walking nodes instead of bisecting; value-bisect \(k\)-th with
     a huge upper bound (\(\coderef{c19-mst-count}\), \(\secref{c19-mst}\))};
  \draw[flow] (verdict) -- (wa);
  \draw[flow] (verdict) -- (mle);
  \draw[flow] (wa) -- (wa1);
  \draw[flow] (wa) -- (wa2);
  \draw[flow] (mle) -- (mle1);
  \draw[flow] (mle) -- (mle2);
\end{tikzpicture}
\caption*{WA: a subtracted root one index off, or an overlap the merge was never idempotent
enough to allow. TLE or MLE: a pool sized for one tree instead of every copy, or a count that
walked instead of bisected.}
\end{figure}


\chaptersection{Quick Reference Card}

\begin{quickref}[Chapter 21 at a glance]
\textbf{Four static / persistent engines, four different ``store more.'' }

\medskip
\textbf{Merge-sort tree (21.1).} A segment tree where each node also stores its interval's
\emph{sorted list} (\(\defref{c19-mst}\)). Build \(\bigO{n\log n}\) by merging on the way up;
count \(\le x\) in \([l,r]\) is one bisection per canonical node, \(\bigO{\log^2 n}\)
(\(\thmref{c19-mstcount}\), \(\coderef{c19-mst-count}\)); \(k\)-th is that count
binary-searched over the \emph{value} axis, \(\bigO{\log^3 n}\).

\medskip
\textbf{Path copying (21.2).} A new version shares every untouched subtree and copies only
the \(\bigO{\log n}\) root-to-leaf spine (\(\defref{c19-pathcopy}\)); keep one root per
version, allocate from a pool sized \(n+q\log n\), never \(4n\) (\(\thmref{c19-memory}\),
\(\coderef{c19-pseg}\)). A lazy tag lives in no single version, so persistence and Ch.~20
refuse to share (\(\secref{c15-strains}\)).

\medskip
\textbf{Chairman / prefix versions (21.3).} \(\mathit{root}_i\) is the value-axis frequency
tree of \(A[1..i]\) (\(\defref{c19-roots}\)); count \(\le x\) in \([l,r]\) is
\(\mathit{root}_r-\mathit{root}_{l-1}\), and the \(k\)-th is a joint two-root descent in
\(\bigO{\log V}\) (\(\thmref{c19-kth}\), \(\coderef{c19-kth}\)). Compress when values are huge.

\medskip
\textbf{Roots on a tree (21.4).} One frequency root per DFS node stores the root-to-\(v\) path
multiset (\(\defref{c19-rootpath}\)); a path \(u\text{--}v\) is a \emph{four}-root combination
\(u,v,\ell,\mathrm{parent}(\ell)\) (\(\propref{c19-fourroots}\), \(\coderef{c19-pathkth}\));
\(\mathrm{lca}\) is Ch.~11's (\(\chref{c11}\)).

\medskip
\textbf{Sparse table (21.5).} \(\mathit{st}[j][i]\) merges the \(2^j\) block from \(i\);
\(\bigO{n\log n}\) build, \(\bigO{1}\) query by \emph{two overlapping} blocks -- valid only for
idempotent merges \(\min/\max/\gcd\) (\(\defref{c19-spacetable}\), \(\thmref{c19-oone}\),
\(\coderef{c19-sparse}\)). Sum uses prefix differences instead (Ch.~18, \(\chref{c14}\)).

\medskip
\textbf{Translation.} ``count / \(k\)-th by \emph{value} in a static range'' \(\to\) merge-sort
tree or chairman. ``read a \emph{past version}'' \(\to\) path copying. ``\(k\)-th on a tree
\emph{path}'' \(\to\) four roots + LCA. ``static \(\min/\max/\gcd\), never updates'' \(\to\)
sparse table. ``modify a whole live range'' \(\to\) lazy (Ch.~20), not here.
\end{quickref}

\begin{lstlisting}[caption={Chapter 21 kit: merge-sort tree (count by value) + persistent point-update/range-sum tree on a node pool + chairman prefix roots (two-root kth) + sparse table, all four in one place}, label={lst:c19-kit}]
// ---- (A) MERGE-SORT TREE: count values <= x in a STATIC range ----
vector<int> ms[4*MX];                       // one SORTED list per node
void msbuild(int v,int l,int r,int a[]){
    if(l==r){ ms[v].push_back(a[l]); return; }
    int m=(l+r)>>1; msbuild(2*v,l,m,a); msbuild(2*v+1,m+1,r,a);
    merge(ms[2*v].begin(),ms[2*v].end(),            // the MERGE in merge-sort
          ms[2*v+1].begin(),ms[2*v+1].end(), back_inserter(ms[v]));
}
int mscount(int v,int l,int r,int ql,int qr,int x){ // # elements <= x in [ql,qr]
    if(qr<l || r<ql) return 0;
    if(ql<=l && r<=qr) return upper_bound(ms[v].begin(),ms[v].end(),x)-ms[v].begin();
    int m=(l+r)>>1;                                  // one bisection per CANONICAL node
    return mscount(2*v,l,m,ql,qr,x)+mscount(2*v+1,m+1,r,ql,qr,x);
}

// ---- (B) PERSISTENT point-update / range-sum tree, ONE node pool ----
struct P{ long long sum; int l, r; } T[POOL];        // T[0] = null
int pk=1; int pnew(long long s,int l,int r){ T[pk]={s,l,r}; return pk++; }
int pbuild(int lo,int hi,int a[]){
    if(lo==hi) return pnew(a[lo],0,0);
    int m=(lo+hi)>>1, L=pbuild(lo,m,a), R=pbuild(m+1,hi,a);
    return pnew(T[L].sum+T[R].sum,L,R);
}
int pupdate(int p,int lo,int hi,int i,long long val){   // returns the NEW root
    if(lo==hi) return pnew(val,0,0);                    // copy the LEAF, do not touch p
    int m=(lo+hi)>>1, L=T[p].l, R=T[p].r;               // share the untouched child
    if(i<=m) L=pupdate(L,lo,m,i,val); else R=pupdate(R,m+1,hi,i,val);
    return pnew(T[L].sum+T[R].sum,L,R);                 // copy this spine node
}
long long pquery(int p,int lo,int hi,int ql,int qr){    // read ANY version root p
    if(qr<lo || hi<ql) return 0;
    if(ql<=lo && hi<=qr) return T[p].sum;
    int m=(lo+hi)>>1;
    return pquery(T[p].l,lo,m,ql,qr)+pquery(T[p].r,m+1,hi,ql,qr);
}

// ---- (C) CHAIRMAN / prefix-version tree: k-th smallest of a[l..r] ----
struct C{ int cnt,l,r; } S[POOL]; int sk=1;             // S[0] = null, cnt 0
int cnew(int c,int l,int r){ S[sk]={c,l,r}; return sk++; }
int cinsert(int p,int lo,int hi,int pos){               // one path copy, +1 on the spine
    int cur=cnew(S[p].cnt+1,S[p].l,S[p].r);
    if(lo!=hi){ int m=(lo+hi)>>1;
        if(pos<=m) S[cur].l=cinsert(S[p].l,lo,m,pos);
        else       S[cur].r=cinsert(S[p].r,m+1,hi,pos); }
    return cur;                                          // ROOT[i]=cinsert(ROOT[i-1],1,V,rk[i])
}
int ckth(int u,int v,int lo,int hi,int k){              // u=ROOT[l-1], v=ROOT[r]
    if(lo==hi) return lo; int m=(lo+hi)>>1;
    int left = S[S[v].l].cnt - S[S[u].l].cnt;           // DIFFERENCE of the two left children
    if(k<=left) return ckth(S[u].l,S[v].l,lo,m,k);
    return ckth(S[u].r,S[v].r,m+1,hi,k-left);
}

// ---- (D) SPARSE TABLE: O(1) idempotent RMQ (min shown; max/gcd/and/or same) ----
int st[20][MX], lg2v[MX];
void sbuild(int a[],int n){
    for(int i=1;i<=n;i++) st[0][i]=a[i];
    lg2v[1]=0; for(int i=2;i<=n;i++) lg2v[i]=lg2v[i/2]+1;
    for(int j=1;(1<<j)<=n;j++)
        for(int i=1;i+(1<<j)-1<=n;i++)
            st[j][i]=min(st[j-1][i], st[j-1][i+(1<<(j-1))]);   // two half blocks
}
int squery(int l,int r){ int j=lg2v[r-l+1];
    return min(st[j][l], st[j][r-(1<<j)+1]); }        // two OVERLAPPING blocks: min/max/gcd ONLY
\end{lstlisting}

\chaptersection{Write-From-Memory Drills}

Five signature skeletons, each lifted from a listing already in the chapter. Cover
the marked line, reproduce from memory, then check the answer one line below.

\begin{lstlisting}[caption={Drill 1 -- merge-sort tree range count: bisect each canonical node (from lst:c19-mst-count)}]
int mscount(int v,int l,int r,int ql,int qr,int x){
    if(qr<l || r<ql) return 0;
    if(ql<=l && r<=qr)
        return upper_bound(ms[v].begin(), ms[v].end(), x) - ms[v].begin();
    int m=(l+r)/2;                                    // partial: split, do NOT bisect here
    return mscount(2*v,l,m,ql,qr,x)+mscount(2*v+1,m+1,r,ql,qr,x);
}
\end{lstlisting}
\noindent{\scriptsize Key: the \code{upper\_bound} happens only on a fully covered node -- the
\(O(\log^2 n)\) is \(O(\log n)\) canonical nodes times one bisection (\(\thmref{c19-mstcount}\)).}

\begin{lstlisting}[caption={Drill 2 -- persistent update: copy the spine and the leaf, share the rest (from lst:c19-pseg)}]
int pupdate(int p,int lo,int hi,int i,long long val){
    if(lo==hi) return pnew(val,0,0);             // 1. COPY the leaf (never write through p)
    int m=(lo+hi)/2, L=T[p].l, R=T[p].r;         // 2. start from p's children (shared)
    if(i<=m) L=pupdate(L,lo,m,i,val); else R=pupdate(R,m+1,hi,i,val);
    return pnew(T[L].sum+T[R].sum,L,R);          // 3. COPY this node; old p untouched
}
\end{lstlisting}
\noindent{\scriptsize Key: exactly \(\bigO{\log n}\) nodes are cloned -- the root-to-leaf path;
every other node is shared with version \(p\) (\(\defref{c19-pathcopy}\), \(\thmref{c19-memory}\)).}

\begin{lstlisting}[caption={Drill 3 -- chairman insert: one path copy, count +1 down the spine (from lst:c19-chairman)}]
int cinsert(int p,int lo,int hi,int pos){
    int cur = cnew(S[p].cnt+1, S[p].l, S[p].r);  // clone node, bump the frequency
    if(lo!=hi){ int m=(lo+hi)/2;
        if(pos<=m) S[cur].l = cinsert(S[p].l,lo,m,pos);
        else       S[cur].r = cinsert(S[p].r,m+1,hi,pos); }
    return cur;
}   // ROOT[i] = cinsert(ROOT[i-1], 1, V, rank(A[i]))
\end{lstlisting}
\noindent{\scriptsize Key: \(\mathit{root}_i\) is a full value histogram, yet only
\(\bigO{\log V}\) nodes are new; the rest is \(\mathit{root}_{i-1}\) reused (\(\defref{c19-roots}\)).}

\begin{lstlisting}[caption={Drill 4 -- chairman k-th: the two-root left-count difference decides each fork (from lst:c19-kth)}]
int ckth(int u,int v,int lo,int hi,int k){     // u=ROOT[l-1], v=ROOT[r]
    if(lo==hi) return lo;
    int m=(lo+hi)/2;
    int left = S[S[v].l].cnt - S[S[u].l].cnt;  // # of a[l..r] whose rank is in [lo,m]
    if(k<=left) return ckth(S[u].l,S[v].l,lo,m,k);
    return     ckth(S[u].r,S[v].r,m+1,hi,k-left);
}
\end{lstlisting}
\noindent{\scriptsize Key: descend BOTH roots in lockstep; \(u\) is \(\mathbf{root}_{l-1}\), not
\(\mathbf{root}_l\) -- the classic off-by-one (\(\thmref{c19-kth}\)).}

\begin{lstlisting}[caption={Drill 5 -- sparse table O(1) query: two overlapping blocks (from lst:c19-sparse)}]
int squery(int l,int r){
    int j = lg2v[r-l+1];                        // largest 2^j not exceeding the length
    return min(st[j][l], st[j][r-(1<<j)+1]);    // left-anchored + right-anchored, they OVERLAP
}
\end{lstlisting}
\noindent{\scriptsize Key: overlap is free only because \(\min(x,x)=x\); swap in \code{sum} and
the intersection is double-counted (\(\thmref{c19-oone}\)).}

\chaptersection{Complete Problem Checklist}

Every problem this chapter offers -- 4 fully solved in-section, worked entries, and one
Summit. Tick \emph{only} after an unaided solve; a hinted attempt counts as half.

\begin{center}
\footnotesize
\setlength{\tabcolsep}{4pt}
\renewcommand{\arraystretch}{1.25}
\begin{tabular}{@{}c p{0.18\linewidth} l p{0.19\linewidth} >{\raggedright\arraybackslash}p{0.23\linewidth} c@{}}
\toprule
\textbf{\#} & \textbf{Problem} & \textbf{Source} & \textbf{Key technique} & \textbf{What to look for} & \textbf{Done?} \\
\midrule
\multicolumn{6}{l}{\textbf{Section 21.1 -- the merge-sort tree}} \\
W21.1 & a count \(\le 3\) over \([2..9]\) on \(A\) & worked & \lv{3} sorted list per node, bisect & stop at canonical nodes; \code{upper\_bound} there only & \(\square\) \\
W21.1b & K-Query & SPOJ KQUERY & \lv{3} count \(>k\) in a range & \(\#\le k\) subtracted from length; \code{int} vs \code{long long} & \(\square\) \\
\midrule
\multicolumn{6}{l}{\textbf{Section 21.2 -- persistence by path copying}} \\
W21.2 & \(A[3]=9\) from \(v_0\) to \(v_1\) on \(A\) & worked & \lv{4} copy exactly 5 of 31 nodes & the leaf is copied; untouched children shared; pool sized \(n{+}q\log n\) & \(\square\) \\
W21.2b & Persistent Bookcase & CF 707D & \lv{4} version tree + DFS rollback & the version graph is a tree, not a line; \code{bitset} per shelf & \(\square\) \\
\midrule
\multicolumn{6}{l}{\textbf{Section 21.3 -- chairman / prefix-version trees}} \\
W21.3 & roots \(0..16\) and kth\([3..9](k{=}4)\) on \(A\) & worked & \lv{4} two-root joint descent & root\({}_{l-1}\) not root\({}_{l}\); compress the value axis & \(\square\) \\
W21.3b & K-th Number & SPOJ MKTHNUM & \lv{4} prefix roots on the value axis & distinct values; \(O(\log n)\) descend, not \(O(\log^3 n)\) & \(\square\) \\
\midrule
\multicolumn{6}{l}{\textbf{Section 21.4 -- persistent roots on a tree}} \\
W21.4 & count \(\le 3\) on path \(3\text{--}8\) of \(T\) & worked & \lv{5} four roots \(u,v,\ell,\mathrm{parent}(\ell)\) & extra root is \(\mathrm{parent}(\ell)\); LCA from Ch.~11 & \(\square\) \\
W21.4b & Count on a tree & SPOJ COT & \lv{5} per-node chairman + k-th on path & \(k\)-th over the path multiset; \(N\log N\) pool & \(\square\) \\
\midrule
\multicolumn{6}{l}{\textbf{Section 21.5 -- sparse table and the engine map}} \\
W21.5 & the full MIN grid on \(A\), query \([3..9]\) & worked & \lv{3} two overlapping blocks & \(j=\lfloor\log_2 L\rfloor\); \code{lg[1]=0}; idempotent only & \(\square\) \\
\midrule
\multicolumn{6}{l}{\textbf{Summit (chapter end)}} \\
W21 & Destiny & CF 840D & \lv{5} merge-sort tree + sampling for \(>\!\mathit{len}/k\) & at most \(k{-}1\) heavy values; smallest qualifying \(x\) & \(\square\) \\
\bottomrule
\end{tabular}
\end{center}

\chaptersection{Recommended Study Order}

Five focused days. Day~1 internalises the merge-sort tree and one count trace; day~2 the
path-copying spine and a version root; day~3 chairman two-root counting and the \(k\)-th
descent; day~4 the four-root path combine; day~5 the sparse table, the engine map, and the
Summit.

\begin{center}
\small
\renewcommand{\arraystretch}{1.25}
\begin{tabular}{@{}llp{0.55\linewidth}@{}}
\toprule
\textbf{Day} & \textbf{Focus} & \textbf{Tasks} \\
\midrule
1 & Read + do 21.1 & re-trace the merge-sort tree over \(A\) (\figref{c19-mst}); one count \(\le 3\) on \([2..9]\) by per-node bisection; write \coderef{c19-mst-build}/\coderef{c19-mst-count} cold; SPOJ KQUERY \\
2 & Read + do 21.2 & the \(A[3]=9\) path copy: name the 5 cloned nodes of 31 and the shared rest; size the pool; then CF 707D end to end (\(\coderef{c19-707-soln}\)) \\
3 & Read + do 21.3 & build roots \(0..16\) on \(A\), verify \(F_9-F_2\); kth\([3..9](k{=}4)=4\) and kth\([1..8](k{=}5)=3\) by joint descent; SPOJ MKTHNUM (\(\coderef{c19-kth}\)) \\
4 & Read + do 21.4 & plant roots on the shared tree \(T\); count \(\le 3\) on path \(3\text{--}8=4\) and kth\((8,10,k{=}6)=4\) from four roots (\figref{c19-path}); SPOJ COT \\
5 & Do 21.5 + Summit & lay the sparse grid on \(A\), query \([3..9]\), reproduce the \(27\ne26\) sum failure; walk the engine map; then the Summit CF 840D; first 3 mixed \\
\bottomrule
\end{tabular}
\end{center}

\chaptersection{Difficulty Progression}

Tiers are competency gates -- each row names the \emph{skill} the tier unlocks.

\begin{center}
\small
\renewcommand{\arraystretch}{1.3}
\begin{tabular}{@{}lp{0.34\linewidth}p{0.38\linewidth}@{}}
\toprule
\textbf{Tier} & \textbf{Problems} & \textbf{What passing this tier proves} \\
\midrule
\lv{3} \ \(\bigstar\bigstar\bigstar\)
  & SPOJ KQUERY, the sparse grid on \(A\), the merge-sort count trace
  & you build a sorted list per node, bisect only on canonical nodes, and read an \(O(1)\) idempotent RMQ \\
\lv{4} \ \(\bigstar\bigstar\bigstar\bigstar\)
  & CF 707D, SPOJ MKTHNUM, the \(A[3]=9\) copy, kth\([3..9]\)
  & you path-copy exactly the spine, size a node pool, and answer a static range \(k\)-th by a two-root descent \\
\lv{5} \ \(\bigstar\bigstar\bigstar\bigstar\bigstar\)
  & SPOJ COT, the four-root path traces, CF 840D (Summit), mixed at CF~2300
  & you combine four roots at the LCA, and choose correctly among merge-sort / chairman / persistent / sparse for a specification \\
\bottomrule
\end{tabular}
\end{center}

\chaptersection{Reduction Gallery}

Every entry is ``one static/persistent flavour + one twist.'' Each line collapses a problem
family onto a listing already in the chapter.

\begin{itemize}[nosep, leftmargin=1.4em]
  \item ``how many of \(a[l..r]\) are \(\le x\) / \(>k\)'' (SPOJ KQUERY) \(\to\) the
        merge-sort tree: one bisection per canonical node
        (\(\thmref{c19-mstcount}\), \(\coderef{c19-mst-count}\), \(\coderef{c19-kquery-soln}\)).
  \item ``\(k\)-th smallest of a static range'' (SPOJ MKTHNUM) \(\to\) chairman prefix roots
        and a joint two-root descent in \(\bigO{\log n}\)
        (\(\thmref{c19-kth}\), \(\coderef{c19-mkthnum-soln}\), \(\probref{c19-mkthnum}\)).
  \item ``read / write a \emph{past version}, branching in time'' (CF 707D) \(\to\) a version
        \emph{tree} with DFS rollback over the persistent state
        (\(\defref{c19-pathcopy}\), \(\coderef{c19-707-soln}\), \(\probref{c19-707d}\)).
  \item ``\(k\)-th / count on a tree \emph{path}'' (SPOJ COT) \(\to\) one frequency root per
        DFS node combined at four roots through Ch.~11's LCA
        (\(\propref{c19-fourroots}\), \(\coderef{c19-pathkth}\), \(\coderef{c19-cot-soln}\)).
  \item ``static \(\min/\max/\gcd\) of a range, no updates ever'' \(\to\) the sparse table's
        two overlapping blocks (\(\thmref{c19-oone}\), \(\coderef{c19-sparse}\)).
  \item ``\(k\)-th on the root-to-\(v\) path (a prefix, not a two-ended path)'' \(\to\) the
        two-root chairman directly on \(\mathit{root}_v\) -- a special case of the four-root
        gallery line above, the persistent-trie law of Ch.~11 (\(\defref{c08-persistent}\)).
  \item ``\(>\!\mathit{len}/k\) majority-with-threshold in a static range'' (CF 840D) \(\to\)
        sample \(O(1)\) positions and count each candidate by the merge-sort tree, or count
        deterministically with chairman roots (\(\secref{c19-mst}\), \(\secref{c19-chairman}\)).
\end{itemize}


\chaptersection{Mixed Problem Set}

Ten entries, all live-verified against the current problem sets and absent from the rest of
the book (the roster's owned IDs -- CF 459D / 522D / 1000F and every Ch.~19/20/11 link -- are
excluded; the first three re-appear only as prose re-treats in Related Topics). None says
``use a merge-sort tree'' or ``use persistence'' -- recognise the engine from \emph{what is
static, what axis the query lives on, and whether history must survive}. Each hint names its
reduction and the Chapter~21 listing it leans on, ordered ascending by effective difficulty.

\begin{mixedset}
  \practice{Count Number of Teams}{LC 1395}
    {https://leetcode.com/problems/count-number-of-teams/}{LC Medium L3}
    {Count index-triplets \(i<j<k\) whose ratings are monotone. Fix the middle \(j\) and the
     answer is (smaller to the left)\(\times\)(larger to the right) plus its mirror -- each
     factor is a ``how many of \(a[\ell..j\!-\!1]\) are \(<x\)?'' query, the count-by-value
     primitive of \(\coderef{c19-mst-count}\) (\(\secref{c19-mst}\)). The \(n\le1000\) bound
     lets a brute scan pass, but the tree turns each half-count into a bisection and is the
     general engine when \(n\) grows (\(\thmref{c19-mstcount}\)).}
  \practice{Frequent Values}{SPOJ FREQUENT}
    {https://www.spoj.com/problems/FREQUENT/}{SPOJ L4}
    {Non-decreasing array, answer the maximum frequency of any value inside \([i,j]\). The
     sorted input means one value occupies a \emph{contiguous run}, so the interior run is the
     interval length and only the two boundary runs need trimming -- a bisection to find each
     value's extent, the same \code{lower/upper\_bound} move that \(\coderef{c19-mst-count}\)
     runs per canonical node (\(\defref{c19-mst}\)). It is the count-occurrences-of-a-value
     question with the merge made free by sortedness (\(\secref{c19-mst}\)).}
  \practice{D-query}{SPOJ DQUERY}
    {https://www.spoj.com/problems/DQUERY/}{SPOJ L4}
    {Count distinct values in \([l,r]\). Read the array as growing prefixes and keep, per
     value, only its \emph{latest} position active: the distinct count in \([l,r]\) is the
     number of active positions \(\ge l\), a two-root difference over the prefix versions --
     exactly the subtraction of \(\thmref{c19-kth}\) (\(\defref{c19-roots}\)). The offline
     right-end sweep with a BIT (Ch.~18, \chref{c14}) is the twin; the persistent form answers
     any \(l,r\) with no reprocessing (\(\coderef{c19-chairman}\)).}
  \practice{K-Query Online}{SPOJ KQUERYO}
    {https://www.spoj.com/problems/KQUERYO/}{SPOJ L4}
    {Count elements \(>k\) in \([a,b]\), but each query is XOR-ed with the previous answer so
     you must answer \emph{online} -- an offline sort-by-right-end sweep can no longer reorder
     them. That is precisely the merge-sort tree's job: bisect the canonical nodes for
     \(\le k\) and subtract from the length (\(\coderef{c19-mst-count}\), \(\thmref{c19-mstcount}\)).
     It is KQUERY (\(\probref{c19-kquery}\)) with the online constraint that kills the BIT
     route and forces a static structure (\(\secref{c19-mst}\)).}
  \practice{Cut and Stick}{CF 1514D}
    {https://codeforces.com/problemset/problem/1514/D}{CF 2000 L4}
    {Can \([l,r]\) be split into at most \(k\) parts so every part is reorderable into a
     palindrome? The obstruction is the values with \emph{odd} frequency, and a value with the
     smallest total frequency can be made even cheaply -- so you need the minimum frequency of
     any value present in \([l,r]\). Binary-search the frequency threshold and, for a candidate,
     count its occurrences by the merge-sort tree's per-node bisection
     (\(\coderef{c19-mst-count}\), \(\secref{c19-mst}\)); the value axis is probed exactly as in
     the \(O(\log^3 n)\) \(k\)-th pipeline (\(\thmref{c19-mstcount}\)).}
  \practice{Max Value of Equation}{LC 1499}
    {https://leetcode.com/problems/max-value-of-equation/}{LC Hard L4}
    {Maximise \(y_i+y_j+|x_i-x_j|\) over pairs within \(|x_i-x_j|\le k\); with \(x\) sorted this
     becomes \((y_i-x_i)+(y_j+x_j)\) with the left partner inside a sliding index window. The
     window maximum of \(y-x\) is an idempotent \(\max\) range query -- either a monotone deque
     or a sparse-table read of two overlapping blocks (\(\thmref{c19-oone}\),
     \(\coderef{c19-sparse}\)). Recognise the ``\(\max\) over a fixed array, no updates'' tell
     that hands it to \(\secref{c19-sparse}\) rather than a mutable tree.}
  \practice{Number of Pairs Satisfying Inequality}{LC 2426}
    {https://leetcode.com/problems/number-of-pairs-satisfying-inequality/}{LC Hard L4}
    {Count \(i<j\) with \(nums1[i]-nums1[j]\le nums2[j]-nums2[i]+\text{diff}\); rearrange to
     \(f(i)\le f(j)+\text{diff}\) where \(f(t)=nums1[t]-nums2[t]\), then for each \(j\) count
     earlier \(i\) with \(f(i)\le f(j)+\text{diff}\) -- a running order-statistic count. Sweep
     \(j\) left to right, insert \(f(i)\) as you go, and query ``how many inserted are \(\le x\)''
     by a bisection on the prefix-sorted structure (\(\coderef{c19-mst-count}\)); the
     two-root prefix view of \(\secref{c19-chairman}\) is the persistent spelling of the same
     count (\(\defref{c19-roots}\)).}
  \practice{Count Good Triplets in an Array}{LC 2179}
    {https://leetcode.com/problems/count-good-triplets-in-an-array/}{LC Hard L5}
    {Two permutations; count value-triplets appearing in increasing position order in
     \emph{both}. Fix the middle value \(v\): the count is (smaller-and-earlier to the
     left)\(\times\)(smaller-and-later to the right), two \emph{order-statistic} counts over the
     position array. A Fenwick gives them offline; the persistent spelling keeps a prefix
     frequency root per index and reads ``how many positions \(<p\) hold a value in a
     range'' as a two-root difference (\(\thmref{c19-kth}\), \(\coderef{c19-chairman}\)) -- the
     count\(\times\)value geometry of \(\secref{c19-mst}\) at \(n\le10^5\).}
  \practice{A Simple Task}{CF 558E}
    {https://codeforces.com/problemset/problem/558/E}{CF 2300 L5}
    {Sort a letter substring ascending or descending, and query the letter-sum of a range, many
     times. Keep at every node a 26-count histogram -- a per-node \emph{sorted multiset} over a
     fixed tiny alphabet, the merge-sort-tree idea of \(\defref{c19-mst}\) with the sorted list
     replaced by counts (\(\coderef{c19-mst-build}\)). A range re-sort writes each letter's count
     into a contiguous stretch. Because the array is now \emph{mutated}, contrast with the static
     tree here: the rewrite itself is Ch.~20's range-assign lazy tag (\chref{c15}), the exact
     boundary \(\secref{c19-copy}\) draws between persistence and mutability.}
  \practice{Optimal Insertion}{CF 1601C}
    {https://codeforces.com/problemset/problem/1601/C}{CF 2300 L5}
    {Pieces of a second array must be inserted into a fixed first array to make the
     concatenation lexicographically smallest; the choice per piece is \emph{where} to drop it.
     The decision reduces to ``how many elements of a value-range sit inside an index-range'' as
     you sweep candidate positions -- a merge-sort-tree count \(\le x\) on the untouched base
     array (\(\coderef{c19-mst-count}\), \(\thmref{c19-mstcount}\)). The base array is static
     across all queries, which is exactly what licenses the sorted-list-per-node engine rather
     than a live tree (\(\secref{c19-mst}\)).}
\end{mixedset}


\chaptersection{The Summit -- CF 840D Destiny}

\begin{notebox}[Why this is the chapter top rope]
Every other problem here reaches for \emph{one} engine. Destiny forces you to hold two in
your head at once and pick per candidate. Its ask -- the smallest value that occurs in a
\emph{static} range \([l,r]\) strictly more than \((r-l+1)/k\) times -- is a counting-by-value
question (Section~21.1) wearing a threshold costume: there are at most \(k-1\le4\) qualifying
values, so you never need to scan the whole value axis, only to \emph{count} a handful of
candidates. The count is exactly the merge-sort tree's bisection (\(\thmref{c19-mstcount}\));
the ``handful'' is the classical heavy-hitter sampling that leans on \(2\le k\le5\). Swap the
randomised hunt for a deterministic one and you are doing Section~21.3's two-root frequency
(\(\thmref{c19-kth}\)) instead -- the same answer, a different engine. \(\secref{c19-sparse}\)'s
map is what tells you \emph{which}.
\end{notebox}

\solvedproblem{c19-summit}{Destiny}{CF 840D}
  {https://codeforces.com/problemset/problem/840/D}{CF 2500 L5}
  {static-range heavy value: sample \(O(1)\) positions of \([l,r]\), count each candidate's frequency by the merge-sort-tree / chairman bisection, keep the smallest that beats \((r-l+1)/k\)}

\begin{problemstatement}[CF 840D Destiny]
An array \(a_1,\dots,a_n\) of integers, \(1\le a_i\le n\). Then \(q\) queries, each three
integers \(l,r,k\) with \(1\le l\le r\le n\) and \(2\le k\le 5\).

\textbf{Answer for a query} is the minimal value \(x\) such that \(x\) occurs in the interval
\([l,r]\) \emph{strictly more than} \(\dfrac{r-l+1}{k}\) times, or \(-1\) if no such number
exists.

\textbf{Input.} First line \(n\) and \(q\) (\(1\le n,q\le 3\cdot10^{5}\)). Next line the \(n\)
values \(a_i\) (\(1\le a_i\le n\)). Then \(q\) lines of \(l,r,k\) (\(2\le k\le 5\)).

\textbf{Output.} One answer per query, on its own line.

\textbf{Sample 1.} \(n{=}4,q{=}2\), array \(1\,1\,2\,2\); queries \((1,3,2)\to\texttt{1}\),
\((1,4,2)\to\texttt{-1}\).\\
\textbf{Sample 2.} \(n{=}5,q{=}3\), array \(1\,2\,1\,3\,2\); queries \((2,5,3)\to\texttt{2}\),
\((1,2,3)\to\texttt{1}\), \((5,5,2)\to\texttt{2}\).
\end{problemstatement}

\begin{attemptgate}[Attempt it yourself first -- \lv{5}, give yourself 40 minutes]
\(n,q\le3\cdot10^5\): a per-query frequency table over the whole window is \(\bigO{n}\) and
times \(q\) is \(\tle\); a full sweep of the value axis to test each candidate is worse.
\peek{the threshold \(>\!(r-l+1)/k\) is generous -- how many values could possibly clear it?}
\nudge{At most \(k-1\) values can each occur more than \((r-l+1)/k\) times (if \(k\) did,
they would sum past the length). With \(k\le5\) there are \(\le4\) ``heavy'' candidates. You
do not have to \emph{find} them by scanning -- \emph{sample} a few positions from \([l,r]\): a
heavy value occupies \(>1/k\) of the window, so a random pick lands on one with probability
\(>1/k\). For each sampled value just \emph{count} its frequency in \([l,r]\).}
\fullhint{For each query take \(S\approx40\) uniform random positions in \([l,r]\), collect the
distinct sampled values, and count each one's occurrences in \([l,r]\) by bisection -- the
merge-sort tree of \(\secref{c19-mst}\) does it in \(\bigO{\log^2 n}\) (\(\coderef{c19-mst-count}\)),
or store, per value, its sorted list of positions and bisect in \(\bigO{\log n}\) (the same
sorted-list-plus-bisection primitive, transposed onto the value axis). Keep the smallest
sampled value whose count \(>\lfloor(r-l+1)/k\rfloor\); if none, print \(-1\). Deterministically,
\(\secref{c19-chairman}\)'s two roots give the exact frequency of any candidate in
\(\bigO{\log n}\), so the sampling is the only probabilistic step.}
\end{attemptgate}

\solstep{Understanding}
The predicate is about a value's \emph{frequency inside a fixed subarray} -- the count-by-value
query Chapter~21 was built for. The window is static, so the merge-sort tree (\(\defref{c19-mst}\))
or the chairman prefix roots (\(\defref{c19-roots}\)) answer ``how many times does \(v\) appear
in \([l,r]\)'' without any updates. The only subtlety is the \emph{search space}: the answer is
not the \(k\)-th smallest but the smallest value exceeding a \emph{frequency} threshold, and
that threshold is so loose (strictly more than a \(1/k\) slice, \(k\le5\)) that few values can
ever qualify.

\solstep{Key Observations}
\begin{observation}
A value ``heavy'' for \((l,r,k)\) has count \(> (r-l+1)/k\). At most \(k-1\le4\) values can be
heavy at once -- \(k\) disjoint groups each larger than a \(1/k\) slice would already overflow
the window (\(\secref{c19-mst}\)'s counting, applied to a budget argument).
\end{observation}
\begin{observation}
Because a heavy value covers a \(>1/k\) fraction of positions, one uniform random position of
\([l,r]\) hits \emph{some} heavy value with probability \(>1/k\). \(S\) samples miss entirely
with probability \(\le(1-1/k)^{S}\), under \(10^{-8}\) for \(k=5,S=40\) -- negligible over all
\(3\cdot10^5\) queries.
\end{observation}
\begin{observation}
``Count occurrences of \(v\) in \([l,r]\)'' is the \emph{identical} primitive the merge-sort
tree uses one level up: the tree bisects a sorted value-list per node
(\(\thmref{c19-mstcount}\)); equivalently, bisect the sorted position-list of value \(v\). Either
way it is the sorted-list-plus-bisection idea of \(\secref{c19-mst}\), and the chairman two-root
difference of \(\thmref{c19-kth}\) is its deterministic twin.
\end{observation}

\solstep{Think It With Me}
\begin{thinkbox}
\deadend{Build a frequency table of the window per query}
{Recounting \([l,r]\) costs \(\bigO{r-l+1}\); summed over \(q\) that is \(\bigO{nq}\) --
\tle, and no structure is being reused between queries.}
\deadend{Binary-search the answer over the value axis like the range \(k\)-th}
{The predicate ``\(x\) occurs \(>\!(r-l+1)/k\) times'' is \emph{not monotone} in \(x\): a large
\(x\) may be heavy while a smaller one is not, so the external bisection of
\(\secref{c19-mst}\) has nothing to sort on. It returns the smallest \emph{heavy} value, not the
\(k\)-th \emph{element}.}
\aha{Turn the threshold into a \emph{budget}: only \(\le k-1\) values can be heavy, so do not
enumerate -- \emph{sample} the window, count each distinct sampled value with the merge-sort
tree / position bisection, and take the minimum that clears the bar. The counting engine of
this chapter supplies the oracle; the sampling is what the loose threshold makes affordable.}
\end{thinkbox}

\solstep{The Approach Ladder}
\begin{center}\renewcommand{\arraystretch}{1.25}
\begin{tabular}{@{}c p{0.22\linewidth} p{0.28\linewidth} l c@{}}
\toprule
\textbf{\#} & \textbf{Approach} & \textbf{Idea} & \textbf{Time} & \textbf{Verdict}\\
\midrule
1 & \lv{2}\ window frequency table & recount \([l,r]\) each query & \(\bigO{nq}\) & \tle \\
2 & \lv{3}\ value-axis binary search & bisect the answer on \(x\) & -- & \textcolor{warningred}{\ding{55}\ \textbf{WA}} \\
3 & \lv{4}\ chairman, every value & two-root count each \(v\) in \(1..n\) & \(\bigO{qn\log n}\) & \tle \\
4 & \lv{5}\ \textbf{heavy sampling + bisection count} & \(S\) picks, count each & \(\bigO{(n+qS)\log n}\) & \ac \\
\bottomrule
\end{tabular}\end{center}
Rung~4 is the Summit: sampling turns ``which value is heaviest'' into ``count a handful of
candidates,'' and each count is one bisection (\(\secref{c19-mst}\) or \(\secref{c19-chairman}\)).
Rung~3 is the deterministic but over-wide version -- exact, and only affordable because a
heavy value is rare enough to hunt probabilistically.

\solstep{Solution}
Preprocess the array into, per value \(v\), its sorted list of positions (or, equivalently,
build the merge-sort tree). For each query \((l,r,k)\): set \(\mathit{need}=\lfloor(r-l+1)/k\rfloor\);
draw \(S\approx40\) uniform random positions in \([l,r]\); for the distinct sampled values,
count their occurrences in \([l,r]\) by bisecting that value's position list; keep the smallest
sampled value whose count is \(>\!\mathit{need}\). Print it, or \(-1\) if no sampled candidate
clears the bar. Repeat the whole query loop with a fresh seed (or larger \(S\)) if a judge
demands extra certainty; the deterministic chairman count of any \emph{proposed} value replaces
the bisection with no change of shape.

\solstep{Proof of Correctness}
\begin{proofbox}
\textbf{Soundness of a printed value.} When the algorithm outputs \(x\), it is because
\(\mathrm{count}(x,[l,r])>\lfloor(r-l+1)/k\rfloor\), and for integers this is exactly
\(\mathrm{count}(x,[l,r])>(r-l+1)/k\) -- the required predicate. It is the \emph{smallest}
among the qualifying candidates it examined. \textbf{Completeness of the candidate set.}
If some value is heavy, at most \(k-1\) are, and each heavy value holds a \(>1/k\) share of the
window, so \(\Pr[\text{a single sample is heavy}]>1/k\ge1/5\); over \(S\) independent samples
\(\Pr[\text{no heavy value ever sampled}]\le(4/5)^{S}<10^{-7}\), and the minimum over all
sampled heavy values equals the true minimum heavy value with that overwhelming probability.
\textbf{Counting correctness.} \(\mathrm{count}(v,[l,r])\) from the position list is
\(\mathit{upper\_bound}(r)-\mathit{lower\_bound}(l)\), the exact number of indices in
\([l,r]\) carrying value \(v\) -- the same bisection the merge-sort tree performs per canonical
node (\(\thmref{c19-mstcount}\)); the chairman variant \(\mathit{root}_r-\mathit{root}_{l-1}\)
is exact by \(\thmref{c19-kth}\). Hence the printed value is the minimal \(x\) satisfying the
predicate, with failure probability negligible in \(S\).
\end{proofbox}

\solstep{Implementation}
\begin{lstlisting}[caption={CF 840D Summit: sample O(1) positions of the static window, count each candidate's frequency by a bisection on its sorted position list (the merge-sort-tree primitive of 21.1, transposed onto the value axis), keep the smallest heavy value over (r-l+1)/k}, label={lst:c19-summit}]
#include <bits/stdc++.h>
using namespace std;
const int MX = 300005;
int a[MX], n, q;
vector<int> pos[MX];              // pos[v] = sorted indices i with a[i]==v
mt19937 rng(712367821U);          // fixed seed -> reproducible output

// frequency of value v in [l,r]: one bisect pair on v's sorted position list.
// The merge-sort tree (21.1) answers the same count in O(log^2 n) via its per-node
// value-list bisection; this is that sorted-list+bisect idea laid out by VALUE, O(log n).
int freq(int v, int l, int r){
    const vector<int>& p = pos[v];
    return (int)(upper_bound(p.begin(),p.end(),r) - lower_bound(p.begin(),p.end(),l));
}

int main(){
    scanf("%d%d", &n, &q);
    for(int i=1;i<=n;i++){ scanf("%d", &a[i]); pos[a[i]].push_back(i); }
    // pos[v] is already increasing because indices are pushed in order i=1..n.
    while(q--){
        int l,r,k; scanf("%d%d%d", &l, &r, &k);
        int need = (r-l+1)/k;                 // answer must occur STRICTLY more than need times
        int best = INT_MAX;
        uniform_int_distribution<int> pick(l, r);
        for(int t=0; t<40; t++){              // 40 samples: miss prob <= (1-1/k)^40 < 1e-7
            int v = a[pick(rng)];
            if(freq(v, l, r) > need && v < best) best = v;
        }
        printf("%d\n", best==INT_MAX ? -1 : best);
    }
}
// Heavy budget: at most k-1 (<=4) values can exceed (r-l+1)/k, so each heavy value covers
// > 1/k of the window and one random position hits a heavy value with prob > 1/k. Sampling
// the window (not scanning the value axis) is what keeps a 3e5-query problem at O(q*log n).
// Deterministic alternative (21.3): a chairman/prefix-root tree gives freq(v in [l,r]) as a
// two-root difference root_r - root_(l-1) in O(log V), removing the random step at a cost.
\end{lstlisting}

\solstep{Trace on the Official Samples}
\begin{tracebox}[Sample 1: array \(1\,1\,2\,2\); position lists \(pos(1)=\{1,2\}\), \(pos(2)=\{3,4\}\)]
\(\mathbf{(1,3,2)}\): len \(=3\), \(\mathit{need}=\lfloor 3/2\rfloor=1\); need count \(>1\).
Sampled values in \([1,3]\) are \(1\) (at indices 1,2) and \(2\) (at index 3).
\(\mathrm{count}(1,[1,3])=2>1\) \(\Rightarrow\) candidate \(1\);
\(\mathrm{count}(2,[1,3])=1\) not \(>1\). Minimum heavy \(=1\). \(\checkmark\)\\
\(\mathbf{(1,4,2)}\): len \(=4\), \(\mathit{need}=2\); need count \(>2\).
\(\mathrm{count}(1,[1,4])=2\), \(\mathrm{count}(2,[1,4])=2\), neither \(>2\) \(\Rightarrow\) no
heavy value \(\Rightarrow\) \(-1\). \(\checkmark\) Output \texttt{1, -1} = official sample 1. \(\Box\)
\end{tracebox}

\noindent\textbf{Replaying sample 2.} Array \(1\,2\,1\,3\,2\), so
\(pos(1)=\{1,3\}\), \(pos(2)=\{2,5\}\), \(pos(3)=\{4\}\).
\((2,5,3)\): len \(=4\), \(\mathit{need}=\lfloor4/3\rfloor=1\); \(\mathrm{count}(2,[2,5])=2>1\)
is the only heavy value, so the answer is \(2\). \((1,2,3)\): len \(=2\),
\(\mathit{need}=\lfloor2/3\rfloor=0\); every present value has count \(\ge1>0\), and the
smallest present is \(1\), so \(1\). \((5,5,2)\): len \(=1\), \(\mathit{need}=0\); the single
value \(2\) has count \(1>0\), so \(2\). The algorithm returns \texttt{2, 1, 2} -- matching the
judge.

\solstep{Complexity}
\begin{complexitybox}
Preprocessing costs \(\bigO{n}\) to fill the position lists (or \(\bigO{n\log n}\) for the
merge-sort tree, matching \(\thmref{c19-mstcount}\)). Each query draws \(S=40\) samples and
runs \(\le40\) distinct bisections at \(\bigO{\log n}\) each (or \(\bigO{\log^2 n}\) via the
tree), so \(\bigO{S\log n}\) per query and \(\bigO{n + qS\log n}\approx3\cdot10^5 +
3\cdot10^5\cdot40\cdot19\approx2.3\cdot10^{8}\) bisect steps -- \ac within 2.5 s under a
\(\bigO{\log n}\) count, and the tree route stays inside 512 MB. Memory \(\bigO{n}\) for the
lists plus the array. The randomised step errs with probability \(<(1-1/5)^{40}\) per query
(\(<10^{-7}\)), negligible across \(3\cdot10^5\) queries; the chairman deterministic route
(\(\secref{c19-chairman}\)) trades the sampling for an exact \(\bigO{\log V}\) count.
\end{complexitybox}

\begin{hardeningbox}
\begin{itemize}[nosep,leftmargin=1.4em]
  \item \textbf{Strictly more than, and integer floor.} The predicate is
        \(\mathrm{count}>\lfloor(r-l+1)/k\rfloor\); using \(<\) or comparing against
        \((r-l+1)/k\) as a real division flips borderline queries like \((1,4,2)\to-1\).
  \item \textbf{Take the minimum heavy value, not the first hit.} Test \emph{all} sampled
        candidates and keep the smallest that clears the bar -- a single accepted sample can
        miss a smaller heavy value.
  \item \textbf{Only count each distinct sampled value once.} Re-bisecting the same value for a
        repeated sample wastes the constant; dedupe per query.
  \item \textbf{Seed and sample count matter.} \(S=40\) is the accepted margin; a tiny \(S\) or a
        clock-seeded RNG that a judge replays can cost you a rare WA -- prefer a fixed seed and
        a larger \(S\) if the checker is adversarial.
  \item \textbf{\(a_i\le n\) bounds the arrays.} The position lists and (if used) the chairman
        value axis stay within \([1,n]\); no coordinate compression is required here, unlike
        \(\defref{c19-roots}\)' general case.
\end{itemize}
\end{hardeningbox}

\begin{recallbox}[Chapter 21 in ten seconds]
When the data is fixed but the questions multiply, store \emph{more}: a \emph{sorted list} in
every interval node (merge-sort tree -- count \(\le x\) by bisection in \(\bigO{\log^2 n}\),
\(k\)-th by bisecting the value axis; \(\defref{c19-mst}\), \(\thmref{c19-mstcount}\)); more
\emph{nodes} -- path-copy the root-to-leaf spine and share the rest, one root per version off a
node pool (\(\defref{c19-pathcopy}\), \(\thmref{c19-memory}\)); more \emph{prefixes} -- the
chairman's \(\mathit{root}_i\) value histograms, so a range is \(\mathit{root}_r-
\mathit{root}_{l-1}\) and its \(k\)-th is a two-root descent (\(\thmref{c19-kth}\)); roots
\emph{planted on a tree}, four of them per path (\(\propref{c19-fourroots}\)); or more
\emph{lengths} -- the sparse table's overlapping power-of-two blocks, \(O(1)\) for idempotent
merges only (\(\thmref{c19-oone}\)). Summit CF~840D (\(\probref{c19-summit}\)) spends the loose
\(>\!(r-l+1)/k\) threshold as a budget of \(\le k-1\) heavy values and pays for each with one
of those bisections.
\selftest{
  \item What single stored object turns ``how many of \(a[l..r]\) are \(\le x\)'' into
        \(\bigO{\log^2 n}\), and how does the same idea give the range \(k\)-th in \(\bigO{\log^3 n}\)? \hfill\lv{3}
  \item Which \(\bigO{\log n}\) nodes does one persistent point-update copy, and why does a node
        pool beat \(4n\) for sizing memory? \hfill\lv{4}
  \item Why is the subtracted chairman root \(\mathit{root}_{l-1}\) and not \(\mathit{root}_l\), and
        why does a \emph{tree path} need \emph{four} roots? \hfill\lv{4}--\lv{5}
  \item Give an operation the sparse table can answer in \(O(1)\) by overlap and one it cannot, and
        say which engine takes the second. \hfill\lv{3}--\lv{4}
  \item In the Summit, why can at most \(k-1\) values clear the threshold, and which chapter bisection
        counts a sampled candidate? \hfill\lv{5}
}
\end{recallbox}


\chaptersection{Tool Kit Inventory}

The machine assembled through Chapter~21. Earlier chapters' rows are carried forward;
this chapter's are appended. Pointers only.

\begin{center}
\small
\renewcommand{\arraystretch}{1.25}
\begin{tabular}{@{}p{0.22\linewidth} l c p{0.24\linewidth} p{0.17\linewidth}@{}}
\toprule
\textbf{Tool} & \textbf{Since} & \textbf{Lv} & \textbf{Trigger phrase} & \textbf{Where} \\
\midrule
\multicolumn{5}{l}{\itshape Carried forward} \\
Plain segment tree (point update, range fold) & 19 & \lv{3} & ``the read engine the sorted lists ride on'' & \chref{c13} \\
Fenwick / difference BIT & 18 & \lv{3} & ``invertible range count, offline sweep'' & \chref{c14} \\
Binary lifting / \(\mathrm{lca}\) & 11 & \lv{4} & ``the meeting point of a tree path'' & \chref{c11} \\
Euler tour \(\tin/\tout\) flatten & 5 & \lv{3} & ``a subtree is a contiguous range'' & \chref{c24} \\
Persistent trie (path copy by depth) & 11 & \lv{4} & ``version roots on a trie'' & \defref{c08-persistent} \\
Lazy tag (mutable whole range) & 20 & \lv{4} & ``a tag lives in no single version'' & \chref{c15} \\
Heavy-light decomposition & 22 & \lv{5} & ``path updates on a changing tree'' & \chref{c16} \\
Segment-tree beats / Li Chao & 23 & \lv{5} & ``non-idempotent \emph{and} mutable combine'' & \chref{c29} \\
\midrule
\multicolumn{5}{l}{\itshape Chapter 21 -- new this chapter} \\
Merge-sort tree (sorted list per node) & 21.1 & \lv{3} & ``count / \(k\)-th by \emph{value} in a static range'' & \defref{c19-mst}, \coderef{c19-mst-count} \\
Path-copying persistent tree & 21.2 & \lv{4} & ``read a \emph{past version} after updates'' & \defref{c19-pathcopy}, \coderef{c19-pseg} \\
Node pool (allocate copies) & 21.2 & \lv{4} & ``size memory by copies, not \(4n\)'' & \thmref{c19-memory} \\
Chairman / prefix-version roots & 21.3 & \lv{4} & ``\(k\)-th of a static range in \(O(\log n)\)'' & \thmref{c19-kth}, \coderef{c19-kth} \\
Persistent roots on a tree (four-root path) & 21.4 & \lv{5} & ``\(k\)-th / count on path \(u\text{--}v\)'' & \propref{c19-fourroots}, \coderef{c19-pathkth} \\
Sparse table (\(O(1)\) idempotent RMQ) & 21.5 & \lv{3} & ``static \(\min/\max/\gcd\), never updates'' & \defref{c19-spacetable}, \coderef{c19-sparse} \\
The engine map & 21.5 & \lv{5} & ``static-vs-mutable \(\times\) axis \(\times\) operation'' & \secref{c19-sparse} \\
Heavy-value sampling (Summit) & Summit & \lv{5} & ``a loose \(>\!\mathit{len}/k\) threshold is few candidates'' & \coderef{c19-summit} \\
\bottomrule
\end{tabular}
\end{center}

\chaptersection{Related Topics}

\begin{itemize}[nosep, leftmargin=1.4em]
  \item \chref{c13} (Ch.~19, Segment Trees) --- backward: Ch.~19 owns the plain
        build / point-update / range-fold tree and the interval geometry this chapter's
        grids reuse; no build/query is re-taught. Its CF 813E Army Creation sliding-window
        solution stays Ch.~19's -- \(\secref{c19-chairman}\) re-treats it in \emph{prose} as the
        chairman twin and never re-links it.
  \item \chref{c15} (Ch.~20, Lazy Propagation) --- backward: Ch.~20 owns range tags;
        \(\secref{c19-copy}\) shows why a lazy tag and a version cannot coexist (a tag lives in
        no single version), back-citing \(\secref{c15-strains}\) with no lazy code here.
  \item \chref{c14} (Ch.~18, Fenwick Trees) --- backward/contrast: prefix sums are the right
        answer where the sparse table's overlap fails (\(\secref{c19-sparse}\)), and the offline
        right-end sweep is the twin the chairman's two-root subtraction replaces (the DQUERY
        gallery line). Three Ch.~18-owned tasks re-enter \emph{this chapter's} ideas only as
        \textbf{prose re-treats, never re-linked}: CF 459D Pashmak and Parmida (per-value
        frequency, the count-by-value oracle of \(\secref{c19-mst}\)), CF 522D Closest Equals
        (nearest equal pair via the last-occurrence prefix view of \(\secref{c19-chairman}\)),
        and CF 1000F One Occurrence (a value occurring exactly once in \([l,r]\), a chairman
        two-root count) -- all three are \emph{owned by Ch.~18 and linked only there}.
  \item \chref{c08} (Ch.~11, Tries) --- backward: the persistent trie (\(\defref{c08-persistent}\),
        solved LC 1707) already path-copied \emph{depth-many} nodes; this chapter generalises to
        \(\log\)-many on a segment tree and cites the trie rather than rebuilding it.
  \item \chref{c11} (Ch.~11, Binary Lifting / LCA) --- backward: \(\secref{c19-ontrees}\) calls
        \(\mathrm{lca}(u,v)\) to name the two subtracted roots of a four-root path and never
        re-derives the lifting table.
  \item \chref{c24} (Ch.~5, Euler Tour) --- backward: \(\tin/\tout\) ``subtree \(=\) interval'' is
        cited when a path query is contrasted with a subtree query; never re-derived here.
  \item \chref{c16} (Ch.~22, Heavy-Light Decomposition) --- forward: a path that is \emph{updated}
        (repaint, add) leaves persistence behind and needs \(O(\log n)\) chain ranges;
        \(\secref{c19-ontrees}\) forward-cites it. This is the \emph{Next} chapter.
  \item \chref{c29} (Ch.~23, Li Chao / Segment Tree Beats) --- forward: a combine that is both
        non-idempotent and mutated is beats' or Li Chao's, not a sparse table or a static tree.
  \item \chref{c21} (Ch.~27, Mo's Algorithm) --- forward: offline range queries with cheap
        add/remove trade the persistent prefix roots for a reordering of the queries.
  \item \chref{c28} (Ch.~17, Rollback / Partial Persistence) --- forward: undo-based \emph{partial}
        persistence is the cheaper cousin of full path copying when only the latest state plus a
        stack of undos is needed; \(\secref{c19-copy}\) names the contrast with no rollback code.
\end{itemize}

\begin{lastword}[Last Word --- Chapter 21]
\emph{``When the data never changes but the questions keep coming, stop folding one
aggregate and start storing more -- more per node (a sorted list), more nodes (a copied
spine of versions), more prefixes (subtractable value histograms), or more lengths (two
overlapping blocks) -- so that counting by value, selecting the \(k\)-th, reading the past,
and asking a static minimum all cost less than one live walk of the range.''}

\medskip
Chapter~19 edited a point and folded a range; Chapter~20 rewrote a whole range with a tag.
Neither answers ``the \(k\)-th smallest of this \emph{fixed} subarray,'' ``how many values
\(\le x\) here,'' or ``what did the tree look like after the third write.'' This chapter's
answer is always the same gesture in a different costume: the merge-sort tree hides a sorted
list in every interval so a bisection replaces a scan (\(\thmref{c19-mstcount}\)); path
copying shares an entire tree and rebuilds only a root-to-leaf spine so history survives at
\(\bigO{\log n}\) per write (\(\defref{c19-pathcopy}\)); the chairman turns time into prefix
histograms that subtract, buying the range \(k\)-th in \(\bigO{\log V}\)
(\(\thmref{c19-kth}\)); four such roots meet at an LCA to read a tree path
(\(\propref{c19-fourroots}\)); and the sparse table keeps every power-of-two window so an
idempotent minimum is two reads (\(\thmref{c19-oone}\)). The Summit CF~840D
(\(\probref{c19-summit}\)) spends a loose threshold as a budget of heavy candidates and pays
for each with one of those bisections.

\medskip
Carry three reflexes forward:
\begin{enumerate}[nosep, leftmargin=1.6em]
  \item \emph{Which axis, and does it move?} A query about \emph{values inside an index range}
        wants a value-sorted structure (merge-sort tree, chairman); a query about a
        \emph{past version} wants path copying; a \emph{pure static idempotent} range wants a
        sparse table; the moment the range is \emph{mutated live}, none of these is right and
        Ch.~20's tag is (\(\secref{c19-sparse}\)'s engine map).
  \item \emph{Subtract the right root.} A range is \(\mathit{root}_r-\mathit{root}_{l-1}\), a
        tree path is four roots and not three -- both off-by-ones are silent WAs
        (\(\thmref{c19-kth}\), \(\propref{c19-fourroots}\)).
  \item \emph{Overlap needs idempotence.} The \(O(1)\) sparse query is legal only because
        re-merging the intersection changes nothing; a sum there is a double-count
        (\(\thmref{c19-oone}\)).
\end{enumerate}

\medskip
And three habits, drilled until boring: budget the node pool for \emph{copies}
(\(n+q\log n\)), never a lone \(4n\); coordinate-compress the chairman's value axis before any
insert when the values are large; and, when a threshold like ``more than a \(1/k\) slice''
appears, ask \emph{how many} answers can qualify before you choose to enumerate or sample. The
mutable point-update tree you built in Chapter~19 is now shadowed by a sorted list, a copied
spine, or a stack of roots -- the same skeleton, storing more, answering the questions one live
aggregate never could.

\medskip
\textbf{Next:} \chref{c16} (Heavy-Light Decomposition) takes the tree-path queries this
chapter answered only for \emph{static} data and makes them \emph{mutable}: it cuts a path
into \(O(\log n)\) contiguous chains so the lazy range updates of Chapter~20 -- the tags
persistence refused to share -- can finally run on a tree, from root to node and node to node.
\end{lastword}
